\documentclass[11pt]{article}
\usepackage[utf8]{inputenc}
\usepackage{amsmath,amssymb}
\usepackage{hyperref}
\title{Robust Atmospheric River Prediction under Distribution Shift}
\author{Anonymous}
\date{}
\begin{document}
\maketitle

\begin{abstract}
This paper studies Atmospheric River Prediction. Grounded in a real, citable literature \cite{ref1,ref2,ref3,ref4,ref5,ref6,ref7,ref8},
we propose a focused method and report \textbf{illustrative} experiments.
All references below are real and verifiable (OpenAlex/arXiv); the reported
numbers are simulated to illustrate intended behaviour.
\end{abstract}

\section{Introduction}
Atmospheric River Prediction is an active research area. This work targets a concrete gap identified from
the recent literature and contributes a method together with an evaluation protocol.

\section{Related Work (real literature)}
The following works, retrieved from public bibliographic APIs, frame our study:
\begin{itemize}
\item Mandelbrot et al. (1968) — "Fractional Brownian Motions, Fractional Noises and Applications", SIAM Review (cited 7723).
\item Lefsky et al. (2002) — "Lidar Remote Sensing for Ecosystem Studies", BioScience (cited 1814).
\item Maurer et al. (2002) — "A Long-Term Hydrologically Based Dataset of Land Surface Fluxes and States for the Conterminous United States*", Journal of Climate (cited 1460).
\item Dai et al. (2002) — "Estimates of Freshwater Discharge from Continents: Latitudinal and Seasonal Variations", Journal of Hydrometeorology (cited 1311).
\item Lam et al. (2023) — "Learning skillful medium-range global weather forecasting", Science (cited 1232).
\item Balsamo et al. (2008) — "A Revised Hydrology for the ECMWF Model: Verification from Field Site to Terrestrial Water Storage and Impact in the Integrated Forecast System", Journal of Hydrometeorology (cited 1129).
\end{itemize}

\section{Method}
We model distribution shift explicitly and optimise a worst-case objective over an uncertainty set of plausible test distributions. We instantiate this idea for Atmospheric River Prediction and analyse its cost/quality trade-off.

\section{Experiments (Illustrative)}
\begin{center}
\begin{tabular}{lcc}
System & Primary Metric & Efficiency \\ \hline
Strong baseline & 0.812 & 1.00$\times$ \\
Ablation & 0.828 & 0.92$\times$ \\
\textbf{Ours} & \textbf{0.851} & \textbf{0.71$\times$}
\end{tabular}
\end{center}
\emph{Metrics simulated to illustrate the intended effect; not measured.}

\section{Conclusion}
We connected a real literature to a concrete method for Atmospheric River Prediction. Future work will
validate the illustrative results with real experiments.

\begin{thebibliography}{9}
\bibitem{ref1} Benoît B. Mandelbrot, John W. Van Ness. Fractional Brownian Motions, Fractional Noises and Applications. \emph{SIAM Review}, 1968. DOI:10.1137/1010093
\bibitem{ref2} M. A. Lefsky, Warren B. Cohen, Geoffrey G. Parker et al.. Lidar Remote Sensing for Ecosystem Studies. \emph{BioScience}, 2002. DOI:10.1641/0006-3568(2002)052[0019:lrsfes]2.0.co;2
\bibitem{ref3} Edwin P. Maurer, Andrew W. Wood, J. C. Adam et al.. A Long-Term Hydrologically Based Dataset of Land Surface Fluxes and States for the Conterminous United States*. \emph{Journal of Climate}, 2002. DOI:10.1175/1520-0442(2002)015<3237:althbd>2.0.co;2
\bibitem{ref4} Aiguo Dai, Kevin E. Trenberth. Estimates of Freshwater Discharge from Continents: Latitudinal and Seasonal Variations. \emph{Journal of Hydrometeorology}, 2002. DOI:10.1175/1525-7541(2002)003<0660:eofdfc>2.0.co;2
\bibitem{ref5} Rémi Lam, Álvaro Sánchez‐González, Matthew Willson et al.. Learning skillful medium-range global weather forecasting. \emph{Science}, 2023. DOI:10.1126/science.adi2336
\bibitem{ref6} Gianpaolo Balsamo, Anton Beljaars, Klaus Scipal et al.. A Revised Hydrology for the ECMWF Model: Verification from Field Site to Terrestrial Water Storage and Impact in the Integrated Forecast System. \emph{Journal of Hydrometeorology}, 2008. DOI:10.1175/2008jhm1068.1
\bibitem{ref7} Charles T. Driscoll, Gregory B. Lawrence, Arthur J. Bulger et al.. Acidic Deposition in the Northeastern United States: Sources and Inputs, Ecosystem Effects, and Management Strategies. \emph{BioScience}, 2001. DOI:10.1641/0006-3568(2001)051[0180:aditnu]2.0.co;2
\bibitem{ref8} Eric Kirby, K. X. Whipple. Quantifying differential rock-uplift rates via stream profile analysis. \emph{Geology}, 2001. DOI:10.1130/0091-7613(2001)029<0415:qdrurv>2.0.co;2
\end{thebibliography}
\end{document}
